% Generated draft: requires mathematical and source review.
\documentclass[11pt]{article}
\usepackage[margin=25mm]{geometry}
\usepackage{fontspec}
\setmainfont{Latin Modern Roman}
\usepackage{amsmath,amssymb,mathtools,unicode-math}
\setmathfont{Latin Modern Math}
\usepackage{enumitem,needspace,xurl,hyperref}
\setlength{\parindent}{0pt}
\setlength{\parskip}{5pt}
\setlength{\emergencystretch}{3em}
\newcommand{\N}{\mathbb{N}}
\newcommand{\Z}{\mathbb{Z}}
\newcommand{\Q}{\mathbb{Q}}
\newcommand{\R}{\mathbb{R}}
\newcommand{\C}{\mathbb{C}}
\newcommand{\F}{\mathbb{F}}
\newcommand{\A}{\mathbb{A}}
\newcommand{\PP}{\mathbb{P}}
\newcommand{\E}{\mathbb{E}}
\newcommand{\eps}{\varepsilon}
\newcommand{\dd}{\,\mathrm{d}}
\begin{document}
\section*{UnsolvedMath ID: 30003995}
\textbf{Source problem code:} catalogue entry
\subsection{Definitions and mathematical statement}
For notation, let \(m,n\) be positive integers and let \(A=(a_{ij})\in\R^{m\times n}\). Define the polyhedron
\[
P(A)=\{x\in\R^n:Ax\leq\mathbf 1_m\},
\]
where the inequalities are coordinatewise and \(\mathbf 1_m\) is the all-ones vector. In the Gaussian model specified by the source, the entries \(a_{ij}\) are mutually independent random variables with distribution \(N(0,1)\). The source specifies no relationship between \(m\) and \(n\).

For a nonempty polyhedron \(P\subseteq\R^n\) and \(v\in P\), its outer normal cone at \(v\) is
\[
N_P(v)=\{u\in\R^n:\langle u,x-v\rangle\leq 0\text{ for every }x\in P\}.
\]
If \(P\) is a bounded polytope, its combinatorial diameter is the diameter of its vertex-edge graph: the maximum, over pairs of vertices, of the minimum number of edges in a path joining them. This is a possible interpretation of the source's term ``diameter,'' but the supplied record does not explicitly specify the diameter convention.

\textbf{Source question, with unresolved specifications.} Can one define a quantitative notion of normal cones being ``wide on average'' and derive useful polyhedral-diameter bounds from it, including a diameter bound for the Gaussian model \(P(A)\) above?

The source does not define cone width, the cones over which to average, their weighting, or the desired bound. It also does not specify whether a Gaussian-model bound should hold in expectation, with a stated probability, or under conditioning. Although \(0\) lies in the interior of \(P(A)\), this polyhedron need not be bounded; the source supplies neither a boundedness hypothesis nor a convention for unbounded realizations. Consequently, the question is a research objective requiring these specifications, rather than a fully quantified conjectural inequality.
\subsection{Short English statement}
The question asks whether an average measure of normal-cone width can control polyhedral diameter, particularly for \(P(A)=\{x:Ax\leq\mathbf 1_m\}\) with independent standard Gaussian matrix entries. The averaging rule, diameter convention, dimensional regime, boundedness treatment, and probabilistic form of the bound remain unspecified.
\subsection{Sources}
\begin{itemize}
\item Supplied local UnsolvedMath registry record, permanent ID \(30003995\), source problem code \texttt{catalogue entry}, titled ``Average Normal-Cone Width and Random Polytope Diameters,'' in the Oberwolfach Reports Open Problems collection. No individual proposer is identified; the supplied proposed year is 2018.
\item Supplied source citation: \emph{Combinatorial Optimization (2019)}, Oberwolfach Reports, DOI 10.4171/owr/2018/50, pp.\ 2969--3023. Supplied source URL and literature reference ``Oberwolfach Report 50/2018'': \url{https://doi.org/10.4171/owr/2018/50}. The differing year labels are retained as supplied.
\item The supplied literature assessment is dated exactly \texttt{2026-08-21} and labels the problem open. It states that the extracted entry asks for an unspecified notion while omitting the mathematical object and conclusion, and classifies the unresolved, unclear claim as open. This is supplied dated curation, not an independently verified literature assessment; the background explicitly requests independent verification of time-sensitive status and cited claims. The accompanying statements supplied here do identify normal cones and a Gaussian polyhedron, so the assessment's description does not fully match those statements.
\item The registry labels the statement \texttt{corrected\_verified} and says that its mathematical content was preserved while notation, typography, duplicate clauses, or omitted definitions were repaired. These are supplied editorial assertions, not independent mathematical verification of this draft or its interpretation.
\end{itemize}
\end{document}
